% paper/sections/scaling_law_iter61.tex -- iter 61 ZVF-conditioned identifiability
% Auto-generated by scripts/scaling_law_iter61.py
\paragraph{Iter 61 elevation: ZVF-conditioned saturation-fit identifiability.}
\label{par:scaling-iter61}
\label{sec:scaling-law-iter61}
\label{sec:iter61}
The iter 25 audit established that 0/5 per-trace saturation-rate
estimates are identifiable at this benchmark's noise floor.  The
iter 57 audit sharpened this to 4/5 anchors at the \(\lambda = 10\)
optimiser bound, with Nemotron-120B as the sole identifiable
exception.  What neither audit explains is \emph{why} the
bound-degeneracy concentrates on these anchors.  This iteration
closes that gap by tying the saturation-fit degeneracy to the same
structural axis that Pillar 2 (\textsc{zvf}) measures: the
per-trace concentration of reward mass at the \(R \in \{0, 1\}\)
extremes.

We define a per-step ZVF proxy
\begin{equation}
  \widetilde{\mathrm{ZVF}} \;=\; \mathbb{P}(R = 0) + \mathbb{P}(R = 1),
  \label{eq:zvf-proxy}
\end{equation}
computed on each anchor's reward trace (\tableref{tab:scaling-iter61-zvf}).
This is the step-level analogue of Pillar 2's rollout-level
\(\mathrm{ZVF} = \mathbb{P}(K_x = 0) + \mathbb{P}(K_x = G)\)
that measures within-group collision.  The two scales differ but
both measure the same property: how much of the binary-outcome
probability mass sits at the all-same extremes.

\begin{table}[t]
  \centering
  \small
  \begin{tabular}{lrrrrrr}
    \toprule
    Model & \(N\) (B) & arch & \(\bar R\) & \(P(R=0)\) & \(P(R=1)\) & \(\widetilde{\mathrm{ZVF}}\) \\
    \midrule
    Qwen3-30B-MoE-Inst & 30.0 & moe & 1.000 & 0.000 & 1.000 & 1.000 \\
    Qwen3-235B-MoE & 235.0 & moe & 1.000 & 0.000 & 1.000 & 1.000 \\
    Nemotron-120B & 120.0 & moe & 0.175 & 0.550 & 0.000 & 0.550 \\
    Qwen3.5-4B & 4.0 & dense & 0.817 & 0.000 & 0.433 & 0.433 \\
    Llama-3.1-8B--Inst & 8.0 & dense & 0.869 & 0.000 & 0.433 & 0.433 \\
    gpt-oss-20B & 20.0 & moe & 0.840 & 0.000 & 0.433 & 0.433 \\
    Kimi-K2-Thinking & 1000.0 & moe & 0.850 & 0.000 & 0.400 & 0.400 \\
    Qwen3.5-27B & 27.0 & dense & 0.437 & 0.333 & 0.000 & 0.333 \\
    DeepSeek-V3.1 & 685.0 & moe & 0.844 & 0.000 & 0.300 & 0.300 \\
    Qwen3-8B & 8.0 & dense & 0.285 & 0.067 & 0.000 & 0.067 \\
    Qwen3-32B & 32.0 & dense & 0.250 & 0.000 & 0.000 & 0.000 \\
    Qwen3-30B-MoE & 30.0 & moe & 0.325 & 0.000 & 0.000 & 0.000 \\
    \bottomrule
  \end{tabular}
  \caption{\textbf{Per-step ZVF proxy across the 12-anchor pool.}
    \(\widetilde{\mathrm{ZVF}} = P(R=0) + P(R=1)\) is the
    step-level counterpart of Pillar 2's rollout-level ZVF.
    Crucially, the extreme-mass concentration comes from two
    distinct sources: \emph{ceiling-degenerate} anchors
    (Qwen3.5-4B, Llama-3.1-8B-Instruct, DeepSeek-V3.1, Kimi-K2,
    Qwen3-30B-MoE-Inst, Qwen3-235B-MoE) carry high \(P(R=1)\) but
    near-zero \(P(R=0)\); the
    \emph{floor-degenerate} anchor Nemotron-120B (LatentMoE, 12B active / 120B total, reclassified as \texttt{moe}) is the unique
    model with high \(P(R=0)=0.55\) and \(\bar R = 0.175\).
    Both archetypes concentrate mass at one extreme, but for
    structurally different reasons -- and the saturation curve
    degenerates in both.  \texttt{scripts/scaling\_law\_iter61.py}
    \(\to\) \texttt{scaling\_law\_iter61\_zvf\_proxy.tsv}.}
  \label{tab:scaling-iter61-zvf}
\end{table}

We stratify the seven full traces into HIGH
(\(\widetilde{\mathrm{ZVF}} \geq 0.1\)) and LOW
(\(\widetilde{\mathrm{ZVF}} < 0.1\)) strata and compare the bounded
saturation fits of \tableref{tab:p1-rawtrace}
(\tableref{tab:scaling-iter61-stratum}).  The short probes are excluded:
a two-parameter fit to 3--5 points is not informative about $\lambda$.

\begin{table}[t]
  \centering
  \small
  \begin{tabular}{lrrr}
    \toprule
    Stratum & \(n\) & \(\hat\lambda\) at bound & mean RMSE \\
    \midrule
    HIGH-ZVF & 6 & 4/6 & 0.186 \\
    LOW-ZVF  & 1 & 1/1 & 0.170 \\
    ALL      & 7 & 5/7 & 0.184 \\
    \bottomrule
  \end{tabular}
  \caption{\textbf{Saturation-fit identifiability by step-level ZVF
    stratum, full traces only.}  Bounded fit $R_{\max} \in [0,1]$,
    $\lambda \in [10^{-4}, 10]$.  HIGH: Qwen3.5-4B,
    Llama-3.1-8B-Instruct, gpt-oss-20B, DeepSeek-V3.1, Nemotron-120B,
    Kimi-K2-Thinking; LOW: Qwen3-8B.  The two interior-$\lambda$ fits
    (gpt-oss-20B, Nemotron-120B) are both HIGH-ZVF, and the only LOW
    trace is on the bound, so ZVF stratum does \emph{not} predict
    bound-degeneracy here.  An earlier version of this table fitted all
    twelve anchors, including 3--5-point probes, with
    $R_{\max} \le 1.5$ and reported 7/9 vs.\ 2/3.  Source:
    \texttt{review\_2026-10-18/P1\_scaling/long7\_zvf\_strata.tsv}.}
  \label{tab:scaling-iter61-stratum}
\end{table}

Nemotron-120B sits in the HIGH-ZVF stratum via its \(P(R=0) = 0.55\)
floor path, which is structurally distinct from the ceiling path
(\(P(R=1) \ge 0.3\)) of the other HIGH anchors.  The ceiling anchors put
$\lambda$ on the bound and \(R_{\max}\) at \(\bar R\) (a near-step
transient); Nemotron-120B has an interior $\hat\lambda = 0.99$ with
\(R_{\max} = 0.182\), but a monotone curve cannot describe its
peak-then-decay trace, so that interior estimate is not meaningful
either.

% REVIEW 2026-10-18: the iter-49 jackknife table (scaling_law_iter61_jackknife.tsv)
% was fitted on summary-reconstructed (synthetic) step traces and is removed,
% together with predictions P3/P4 and panel C of the former iter-61 figure.
The cross-pillar ZVF comparison
(\tableref{tab:scaling-iter61-cross}) shows that the per-step
ZVF proxy across the 12 anchors is
\(\widetilde{\mathrm{ZVF}}_{\mathrm{step}} = 0.413\)
versus Pillar 2's rollout-level
\(\mathrm{ZVF}_{\mathrm{rollout}} = 0.158\)
(Qwen3-8B, G=8, 600 problems).  The step-level proxy is
+0.254 units
{\emph{higher}} than the rollout-level ZVF -- the expected
direction, because per-step aggregation cannot de-correlate
within-group successes that the rollout-level grouping sees as
mixed.  The two quantities are computed at different granularities
and on different runs, so this comparison is consistent with, but does
not test, the Pillar-2 reading that the degeneracy is a within-group
contrast problem.

\begin{table}[t]
  \centering
  \small
  \begin{tabular}{lll}
    \toprule
    Source & Metric & Value \\
    \midrule
    Pillar2\_tinker\_G8 & mean\_zvf\_per\_group & 0.1583 \\
    Pillar2\_tinker\_G8 & p\_extreme\_per\_group & 0.1584 \\
    Pillar2\_tinker\_G8 & p\_mixed\_per\_group & 0.8417 \\
    Pillar1\_iter61\_step\_zvf & mean\_zvf\_proxy\_per\_step & 0.4125 \\
    Pillar1\_iter61\_step\_zvf & median\_zvf\_proxy\_per\_step & 0.4167 \\
    Pillar1\_iter61\_step\_zvf & mean\_p\_zero\_per\_step & 0.0792 \\
    Pillar1\_iter61\_step\_zvf & median\_p\_zero\_per\_step & 0.0000 \\
    Prediction & step\_zvf\_over\_rollout\_zvf & 2.6058 \\
    \bottomrule
  \end{tabular}
  \caption{\textbf{Cross-pillar ZVF: rollout-level (Pillar 2) vs
    step-level (Pillar 1 iter 61).}  The two scales agree on
    direction (extreme-mass concentration predicts bound-degeneracy)
    but differ in magnitude because per-step aggregation cannot
    de-correlate within-group successes.  Pillar 2's 15.8\%
    rollout-level ZVF is the per-group estimate at G=8; the iter 61
    step-level proxy is the per-step aggregation across the
    12-anchor pool, which is necessarily higher because each
    step contributes an entire group worth of variance.
    \texttt{scripts/scaling\_law\_iter61.py}
    \(\to\) \texttt{scaling\_law\_iter61\_cross\_pillar.tsv}.}
  \label{tab:scaling-iter61-cross}
\end{table}

\paragraph{Internally recorded predictions.}
Three of the five predictions in \texttt{scaling\_law\_iter61\_predictions.tsv} are retained (P3 and P4 concerned a jackknife of a cross-scale fit computed on synthetic traces and are withdrawn):
\begin{itemize}
  \item \texttt{P1\_high\_stratum\_more\_bound}: \textbf{FAIL} on the recomputed full-trace strata ($4/6 = 0.67$ HIGH vs.\ $1/1$ LOW; the recorded PASS, $0.78 > 0.67$, came from the twelve-anchor fit including 3--5-point probes).
  \item \texttt{P2\_nemotron\_in\_high\_floor}: \textbf{PASS}. Nemotron-120B sits in HIGH-ZVF via $P(R=0)=0.55$ (floor-collapse, not ceiling) ($\text{observed}=\text{True}$, $\text{expected}=\text{True}$).
  \item \texttt{P5\_step\_zvf\_above\_rollout\_zvf}: \textbf{PASS}. Mean per-step ZVF proxy exceeds Pillar 2 rollout-level $\text{mean\_zvf}$ ($\text{observed}=0.4125 > 0.1583$).
\end{itemize}

\paragraph{What iter 61 shows.}
The iter 25/57/21 audits established that the canonical saturation
fit is degenerate on this benchmark.  Iter 61 describes two distinct
ways a trace concentrates reward mass at an extreme -- ceiling-collapse ($P(R=1)$ high, anchors that have
saturated) and floor-collapse ($P(R=0)$ high, the Nemotron-120B
counterexample) -- but on the full traces the stratum does not predict which fits
reach the $\lambda$ bound, so this is a description rather than an
explanation of the degeneracy.  Pillar 2's contrastive-yield
decomposition remains a candidate explanation that these single-seed
step-level summaries cannot test; rollout-level ZVF logged on the same
runs would be needed.

% (iter-61 figure removed: panel C plotted the synthetic-trace jackknife.)
